% !TEX root = ../main.tex
% =====================================================================
%  §4  Twisted N=2 SYM In Mathai-Quillen Form
%  原文：pp. 19--37（PDF 第 20--38 页）；公式 (4.1)--(4.76)
%  说明：正文中 "$s(A):=F^+:=F+*F$"（式 (4.2)）不含 1/2 因子，
%        全文的自对偶/反自对偶约定均以此为准。
% =====================================================================

\section{扭曲 $\mathcal{N}=2$ 超对称杨--米尔斯理论的 Mathai--Quillen 形式}
\label{sec:mq-donaldson}

\begin{keytakeaways}
\item \textbf{物理图像：}把 $\mathcal{N}=2$ 超对称杨--米尔斯理论放在一般（未必自旋的）黎曼
  四流形上，只需额外引入一个满足 $w_2(P_R)=w_2(X)$ 的 $SO(3)$ $R$-对称丛 $P_R$；再把
  $P_R$ 与自对偶标架丛 $P^+$ 等同（$\omega_R=\omega^+$），即完成拓扑扭曲，得到以 $Q$ 为
  BRST 算子的上同调拓扑场论，其路径积分局域化到反自对偶瞬子模空间。
\item \textbf{数学工具：}Mathai--Quillen/CohTFT 的局域化机器（场 $A,H,\phi,\bar\phi;\psi,\chi,\eta$
  与形变复形 $0\to\Omega^0\xrightarrow{\nabla_A}\Omega^1\xrightarrow{\nabla^+_A}\Omega^{2,+}$）、
  等变上同调的嘉当模型与下降（descent）形式、阿蒂亚--希钦--辛格的椭圆复形指标，
  以及唐纳森的 $\mu$ 映射与斜积。
\item \textbf{核心结论：}$Q$-闭可观测量的关联函数在 $\omega_R=\omega^+$ 处与度量无关，
  等于唐纳森多项式（相差整体常数）；虚拟维数 $4h^\vee k-\dim G\,(b_2^+-b_1+1)$ 同时
  就是 $u(1)_R$ 鬼数反常，故只留下 $4\ell+2r=\vdim\mathcal{M}$ 的项；配分函数
  (4.59) 中的 $\Lambda$ 不能简单置 1，它记录了 UV 理论的重整化标度，是 §5 以后把 UV
  理论映到 IR 有效理论的关键。
\end{keytakeaways}

% [Tutor Note: 本节是全篇的技术枢纽：§3 建立的抽象 Mathai--Quillen 机器在这里第一次被
%  真正“接线”到物理的 $\mathcal{N}=2$ SYM 上。阅读时请始终盯住两条线索：(i) 鬼数
%  $=$ 微分形式次数 $=$ $u(1)_R$ 荷这一三重等同；(ii) $Q$ 的标量性（式 (4.16)）如何由
%  拓扑扭曲保证。这两条合起来给出 (4.59) 的选择定则与唐纳森多项式的物理诠释。]

% ---------------------------------------------------------------------
%  §4.1
% ---------------------------------------------------------------------
\subsection{场、方程、对称性}
\label{sec:4-1}

基本数据为：
\begin{enumerate}[label=\arabic*., leftmargin=2.2em]
  \item 紧李群 $G$，其李代数为 $\mathfrak{g}$；
  \item 一个闭、定向的黎曼 $4$-流形 $(X,g_{\mu\nu})$；
  \item 一个主丛 $P\to X$。
\end{enumerate}

我们将把坐标记为 $x^i$ 的场空间 $M$ 换成联络空间，即 $A\in\mathcal{A}:=\Conn(P)$；
把方程丛 $E\to M$ 换成
\begin{equation}
  \mathcal{E}=\mathcal{A}\times\Omega^{2,+}(X,\mathrm{ad}P)
  \tag{4.1}
\end{equation}
而截面取为
\begin{equation}
  s(A):=F^{+}:=F+\ast F .
  \tag{4.2}
\end{equation}
方程丛与截面关于对称群是等变的：
\begin{equation}
  \mathcal{G}=\Aut(P).
  \tag{4.3}
\end{equation}

% [Tutor Note: 原文此处写作“$G$ 是该理论的规范群，而 $\mathcal{G}$ 是规范变换群”，但前
%  一句又已令 $\mathcal{G}=\Aut(P)$。按后文（尤其 §4.2.1 第 4 条）的用法，斜体 $G$ 始终指
%  结构群，花体 $\mathcal{G}$ 才指规范变换群（自同构群）。译文保留原文的字体区分，读者不
%  必把这句话当作两种不同的群。另注意 (4.2) 的约定 $F^+:=F+\ast F$ 不含 $1/2$ 因子，这与
%  文献中常见的 $F^\pm=\tfrac12(F\pm\ast F)$ 不同；本讲义全部符号均以此为准。]

于是 $G$ 是该理论的结构群，而 $\mathcal{G}$ 是规范变换群；局部上可以把 $\mathcal{G}$
的元素看作由映射 $X\to G$ 给出的规范变换。

{\sloppy
这样我们就能开动第 3 节所回顾的上同调场论的机器。按照那里的形式结构，对
$(A,H,\phi,\bar\phi;\psi,\chi,\eta)$ 的路径积分将局域化到瞬子模空间
$\mathcal{M}(P,g)$ 上；下面我们说明后者与物理的超对称杨--米尔斯理论有何关系。\par}

% ---------------------------------------------------------------------
%  §4.2
% ---------------------------------------------------------------------
\subsection{与扭曲 $\mathcal{N}=2$ 场论的关系}
\label{sec:4-2}

下面详细说明四维拓扑扭曲的 $\mathcal{N}=2$ 量子场论（QFT）是什么意思。

在物理上，一个 $\mathcal{N}=2$ 的 QFT 是指在 $3+1$ 维时空、其希尔伯特空间上带有
$\mathcal{N}=2$ 超庞加莱代数的一个酉表示。经 Wick 转动到欧氏符号差后，该超代数的偶部为
\begin{equation}
  \mathfrak{SP}^{0}=\mathbb{R}^{4}\rtimes
  \bigl(\mathfrak{su}(2)_{-}\oplus\mathfrak{su}(2)_{+}\bigr)
  \oplus\mathfrak{su}(2)_{R}\oplus\mathfrak{u}(1)_{R},
  \tag{4.4}
\end{equation}
奇部为
\begin{equation}
  \mathfrak{SP}^{1}=(1;2;2)^{1}\oplus(2;1;2)^{-1}.
  \tag{4.5}
\end{equation}

% [Tutor Note: 原文紧接下文写“$Q^{A}_{\alpha}$ 属于 $(2;1;1)^{-1}$”，与 (4.5) 中第二项的
%  $(2;1;2)^{-1}$ 不符——$R$-对称的指标 $2$ 显然脱落了，$(2;1;1)$ 也不是
%  $\mathfrak{su}(2)_-\oplus\mathfrak{su}(2)_+\oplus\mathfrak{su}(2)_R$ 的合法表示。应以
%  (4.5) 为准，$Q^{A}_{\alpha}$ 属于 $(2;1;2)^{-1}$。译文照录原文的 $(2;1;1)^{-1}$ 以免
%  改动原文，但读者应径自改为 $(2;1;2)^{-1}$。]

$\mathfrak{SP}^{0}$ 的前两个直和项是转动生成元，按单位四元数的左作用与右作用分解；
后两个直和项的下标 “$R$” 表示它们是“$R$-对称性”。一般地，$R$-对称性是不与超荷对易的
整体对称性。

我们把奇部写成了 $\mathfrak{SP}^{0}$ 的表示；上标指 $\mathfrak{u}(1)_R$ 的荷。存在一个
等变映射
\begin{equation}
  \Sym^{2}(\mathfrak{SP}^{1})\to\mathbb{R}^{4}
  \tag{4.6}
\end{equation}
满足超雅可比关系。

物理学家通常把奇生成元记作超荷：$\bar{Q}^{A}_{\dot\alpha}$ 属于 $(1;2;2)^{+1}$，
$Q^{A}_{\alpha}$ 属于 $(2;1;1)^{-1}$。指标 $\alpha,\dot\alpha,A$ 均取 $1,2$，各自在相应
因子上标记二重态。在物理文献中映射 (4.6) 几乎总写成
\begin{equation}
  \begin{aligned}
    \{Q^{A}_{\alpha},\bar{Q}^{B}_{\dot\beta}\}
      &=2\epsilon^{AB}\sigma^{\mu}_{\alpha\dot\beta}P_{\mu},\\
    \{Q^{A}_{\alpha},Q^{B}_{\beta}\}
      &=0 .
  \end{aligned}
  \tag{4.7}
\end{equation}
此外，当 $Q,\bar{Q}$ 是作用在物理希尔伯特空间上的算子时，还有实条件
$(Q^{A}_{\alpha})^{\dagger}=\epsilon_{AB}\bar{Q}^{B}_{\dot\alpha}$。\footnote{在拓扑
扭曲理论中这个实条件将被破坏。事实上可以证明，不存在任何酉的拓扑场论能重现唐纳森
不变量 \cite{ref13}。}

存在 $\mathfrak{su}(2)_{R}$ 对称性正是 “$\mathcal{N}=2$” 这一名称的来源——它意味着有两
组 $\mathcal{N}=1$ 超对称生成元。

在 $\mathcal{N}=2$ 场论中，场按超对称代数的表示变换。有两类对我们至关重要的场表示：
“矢量多重态”与“超多重态”。

% ---------------------------------------------------------------------
%  §4.2.1
% ---------------------------------------------------------------------
\subsubsection{$\mathcal{N}=2$ 矢量多重态}
\label{sec:4-2-1}

矢量多重态的场列于下表。

\begin{center}
\renewcommand{\arraystretch}{1.18}
\begin{tabular}{|c|c|c|}
\hline
 & $\mathfrak{su}(2)_{-}\oplus\mathfrak{su}(2)_{+}\oplus\mathfrak{su}(2)_{R}$
 & $\mathfrak{u}(1)_{R}$\\
\hline
$A_{\mu}$                  & $(2,2,1)$ & $0$\\
\hline
$\bar{\psi}^{A}_{\dot\alpha}$ & $(1,2,2)$ & $-1$\\
\hline
$\psi^{A}_{\alpha}$        & $(2,1,2)$ & $1$\\
\hline
$\phi$                     & $(1,1,1)$ & $2$\\
\hline
$\bar{\phi}$               & $(1,1,1)$ & $-2$\\
\hline
$D$                        & $(1,1,3)$ & $0$\\
\hline
\end{tabular}
\end{center}

% [Tutor Note: 这张表是全书反复引用的“字典”。表中 $D$ 在 $\mathfrak{su}(2)_R$ 下是三重态
%  $(1,1,3)$；扭曲后 $\mathfrak{su}(2)'_+$ 被认同为自对偶洛伦兹群，而 $\mathbf 3$ 正是
%  自对偶二形式 $\Lambda^{2,+}$ 的表示（$\Lambda^{2,+}\cong\Sym^{2}S^{+}$），故此三重态
%  整体成为自对偶二形式 $D^{+}_{\mu\nu}$（见 (4.31)），并非“只有某个分量”是辅助场。
%  $\phi,\bar\phi$ 的 $u(1)_R$ 荷为 $\pm2$（即鬼数 $\pm2$）：扭曲后 $\phi$ 充当嘉当模型中
%  次数为 $2$ 的等变参数而 $\bar\phi$ 充当投影多重态的一半，两者角色全然不同——这正是
%  原文强调“从不使用复共轭关系”的原因。]

\begin{enumerate}[label=\arabic*., leftmargin=2.2em]
  \item 我用的是数学家的 $\mathfrak{su}(2)$ 表示记号，即以维数标记表示。

  \item $\mathfrak{u}(1)_R$ 量子数将对应于“鬼数”，而后者又对应于瞬子模空间上的微分形式
    次数。

  \item 在物理理论中这些场满足实条件，且全部取值于 $\mathfrak{g}$ 的伴随表示。$\phi$
    与 $\bar\phi$ 互为复共轭，故可写 $\phi=B_{1}+\ii B_{2}$、$\bar\phi=B_{1}-\ii B_{2}$，
    其中 $B_{1},B_{2}$ 是取值于 $\mathfrak{g}$ 的标量场。在拓扑扭曲理论中 $\phi$ 与
    $\bar\phi$ 扮演极为不同的角色，从不使用复共轭关系；对 $B_{1},B_{2}$ 的路径积分被
    替换为对 $\phi$ 与 $\bar\phi$ 的“围道积分”。

  \item 上面给出的局部洛伦兹群下的量子数，告诉我们如何把矢量多重态的场解释为一般主
    $G$-丛 $P\to X$（$X$ 为一般四流形）上的对象。为此我们\textbf{必须}引入 $R$-对称
    丛这一额外数据：
    \begin{equation}
      P_{R}\to X .
      \tag{4.8}
    \end{equation}
    $P_{R}$ 是主 $SU(2)$ 丛，是把矢量多重态放到一般四流形 $X$ 上所需的额外数据。于是：
    \begin{enumerate}[label=\arabic*., leftmargin=1.6em]
      \item $A$ 是联络；
      \item $\phi$ 与 $\bar\phi$ 是 $\mathrm{ad}P\otimes\mathbb{C}$ 的截面；
      \item $D$ 是 $\mathrm{ad}P\otimes W$ 的截面，其中 $W\to X$ 是与 $P_{R}$ 相伴的
        实秩 $3$ 向量丛；
      \item $\psi$ 与 $\bar\psi$ 是 $S^{\pm}\otimes S_{R}\otimes\mathrm{ad}(P)$ 的截面，
        其中 $S^{\pm}$ 是 $X$ 的自旋丛，$S_{R}$ 是……
    \end{enumerate}

  \item 在这一课题中，考虑非自旋的四流形 $X$ 相当重要（因为“绝大多数”四流形都不是自旋
    的）。\footnote{关于自旋结构及相关事项的速成介绍见附录 A。}当 $X$ 非自旋时，丛
    $S^{\pm}$ 并不存在。此时我们应取 $R$-对称丛为主 $SO(3)$ 丛，并且必须有
    \begin{equation}
      w_{2}(X)=w_{2}(P_{R}).
      \tag{4.9}
    \end{equation}
    这样 $S^{\pm}\otimes S_{R}$ 就存在，即使各个因子单独并不存在。这是红外理论中必须
    引入自旋$^c$ 结构这一要求的紫外对应物。事实上，我们将把结构群为 $SO(3)$ 的
    $R$-对称丛取为 $X$ 上与自对偶形式丛相伴的主 $SO(3)$ 丛 $P^{+}$，即选定一个同构
    \footnote{〔原稿边注〕同构的选取是否会带来有意思的依赖？}
    \begin{equation}
      P_{R}\cong P^{+}.
      \tag{4.10}
    \end{equation}
\end{enumerate}

% [Tutor Note: 第 4 条末尾“$S_R$ 是……的自旋表示”在原稿中被截断，句子停在
%  “associated to an” 处而没有下文。按上下文，$S_R$ 应是与 $P_R$ 相伴的（主 $SU(2)$ 丛
%  的）自旋表示丛，其存在性正由 (4.9) 的 $w_2$ 条件保证。此处照实译出并不补全，以免
%  引入原文所无的内容。]

% [Tutor Note: 本节末尾转入 §4.2.2。作者在讲义中始终区分“把 $P_R$ 与 $P^+$ 同构”和
%  “令 $\omega_R=\omega^+$”两件事：前者是丛的同构（拓扑数据），后者是该同构下联络的
%  等同（几何数据）。(4.13) 用的是后者。这个区分在 $b_2^+=1$ 时与穿墙现象相关，见 §4.6
%  末尾的提醒。]

% ---------------------------------------------------------------------
%  §4.2.2
% ---------------------------------------------------------------------
\subsubsection{拓扑扭曲}
\label{sec:4-2-2}

威滕 1988 年论文的关键创新之一是拓扑扭曲的概念。思考这种拓扑扭曲有两种有用的方式。

如我们所见，为了把矢量多重态的场解释为一般四流形 $X$ 上一般主 $G$-丛 $P\to X$ 上的
对象，我们必须同时引入一个满足 $w_2(X)=w_2(P_R)$ 的 $SO(3)$ $R$-对称丛 $P_R$。为了写出
作用量、从而写出路径积分，我们在 $X$ 上引入黎曼度量 $g_{\mu\nu}$，并且还必须在 $P_R$
上引入一个联络 $\omega_R$。$X$ 上的黎曼度量给出 $TM$ 上的 Levi-Civita 联络；把联络
$1$-形式分解为（反）自对偶分量 $\omega^{\pm}$ 后即可看出，路径积分将是三个外场的函数：
\begin{equation}
  Z(\omega^{-},\omega^{+},\omega_{R})=\int[\dd A\,\dd\psi\cdots]\,\ee^{S_{\mathrm{phys}}},
  \tag{4.11}
\end{equation}
其中
\begin{equation}
  \begin{aligned}
    S_{\mathrm{phys}}
    &=-\int_{X}\frac{1}{g_{0}^{2}}\tr\Bigl(
        F\wedge\ast F+D\phi\wedge\ast D\phi^{*}
        -\frac{1}{4}[\phi,\phi^{*}]^{2}\,\vol\Bigr)\\
    &\quad+\frac{\theta_{0}}{8\pi^{2}}\int_{X}\tr F\wedge F
      +S_{\mathrm{phys,Fermi}},
  \end{aligned}
  \tag{4.12}
\end{equation}
其中 $\tr(XY)=-\frac{1}{2h^{\vee}}\Tr_{\mathfrak{g}}\ad(X)\ad(Y)$ 是李代数
$\mathfrak{g}$ 上的一个正定形式，$\vol$ 是由度量 $g_{\mu\nu}$ 诱导的体积形式，
$g_{0}^{2}$ 是裸杨--米尔斯耦合，$S_{\mathrm{phys,Fermi}}$ 是含费米子的项。事实上，对
$P_{R}$ 上联络 $\omega_{R}$ 的依赖只出现在费米子项的动能里。

% [Tutor Note: (4.12) 的迹归一化 $\tr(XY)=-\frac{1}{2h^{\vee}}\Tr_{\mathfrak{g}}
%  \ad(X)\ad(Y)$ 中的负号不是笔误：对紧致 $\mathfrak{g}$，$\ad X$ 反称，故
%  $\Tr_{\mathfrak{g}}(\ad X\,\ad X)<0$，加上负号后 $\tr$ 才正定。这个归一化并非无关
%  紧要的约定：正是它使得 $\int_X\tr(F\wedge F)$ 在任意 $X$ 上取整数值、且对适当的
%  $P$ 取遍所有整数（见 (4.17) 之后的说明），从而 $\theta$ 角的周期为 $2\pi$、瞬子数
%  $k$ 为整数（(4.27)）。若换成物理文献中常见的 $\tr_{\mathrm{fund}}$ 归一化，则
%  (4.17) 的 $\tau_{0}$ 定义 (4.18) 中的系数 $4\pi\ii/g_{0}^{2}$ 要相应改变。]

现在按上面的方式选定同构 (4.10)，并把配分函数 (4.11) 限制到满足
\begin{equation}
  \omega^{+}\cong\omega_{R}
  \tag{4.13}
\end{equation}
的外场上（“$\cong$”指在该同构下相等）。于是奇妙的事情发生了：

\begin{center}
\itshape 对 $\omega^{-}$ 与 $\omega^{+}$ 的依赖同时消失！
\end{center}

% [Tutor Note: 这一“奇妙的事情”只在形式层面成立。要理解它为何成立，最方便的是切换到
%  第二种（也是威滕原文 [53] 的）视角：把 $\omega_R=\omega^+$ 看作改变理论对引力场
%  $\omega^{\pm}$ 的耦合，而引力耦合由场的洛伦兹量子数决定；新的耦合等价于选取
%  (4.14)(4.15) 的“扭曲”洛伦兹群。此时原来分别属于 $\mathfrak{su}(2)_+$ 与
%  $\mathfrak{su}(2)_R$ 的指标被合并，$R$-对称的指标变成自旋指标，于是场的几何意义
%  （即它是联络、$1$-形式、自对偶 $2$-形式还是标量）完全由 $(\mathfrak{su}(2)_-,
%  \mathfrak{su}(2)'_+)$ 的表示论读出。]

至少，在形式层面发生的正是这件事。要看出原因，方便的做法是回顾关于拓扑扭曲的第二种
视角（这也是威滕原文 \cite{ref53} 的原始视角）。

令 $\omega_{R}=\omega^{+}$，就是在改变理论对引力场 $\omega^{\pm}$ 的耦合。对引力的耦合
由场的洛伦兹量子数决定，而这一耦合的改变可以描述为选取一个同构的洛伦兹群，其李代数为
\begin{equation}
  \mathfrak{su}(2)_{-}\oplus\mathfrak{su}(2)'_{+},
  \tag{4.14}
\end{equation}
其中
\begin{equation}
  \mathfrak{su}(2)'_{+}:=\mathrm{Diag}\subset
  \mathfrak{su}(2)_{+}\oplus\mathfrak{su}(2)_{R}.
  \tag{4.15}
\end{equation}
其实际后果是：只需把上表中 $\mathfrak{su}(2)_{-}$ 与 $\mathfrak{su}(2)_{R}$ 的相应表示
作张量积，并把它们看作同一个 $\mathfrak{su}(2)$ 的表示，就能读出场的几何诠释。

这种拓扑扭曲的关键动机在于：我们需要一个能充当 BRST 算子的超对称算子。这意味着：
\begin{enumerate}[label=\arabic*., leftmargin=2.2em]
  \item 它的平方必须为零；
  \item 它必须是标量，从而可以在任意黎曼 $4$-流形上定义。
\end{enumerate}
这就促成了上面定义的扭曲，我们取
\begin{equation}
  Q=\delta^{\dot\alpha}_{A}\bar{Q}^{A}_{\dot\alpha}.
  \tag{4.16}
\end{equation}

% [Tutor Note: (4.16) 之所以是标量，正是 (4.14)(4.15) 的“对角”扭曲的直接结果。扭曲前
%  $\bar Q^A_{\dot\alpha}$ 属于 $(1;2;2)^{+1}$：$\dot\alpha$ 是 $\mathfrak{su}(2)_+$ 的
%  右旋二重态指标，$A$ 是 $\mathfrak{su}(2)_R$ 的二重态指标。取对角子代数
%  $\mathfrak{su}(2)'_+=\mathrm{Diag}$ 后，$\mathbf 2\otimes\mathbf 2=\mathbf 1\oplus
%  \mathbf 3$，而用不变张量 $\delta^{\dot\alpha}_{A}$ 缩并恰好取出其中的单态部分，故
%  $Q$ 在扭曲后的洛伦兹群 $\mathfrak{su}(2)_-\oplus\mathfrak{su}(2)'_+$ 下是标量，可以
%  在任意黎曼 $4$-流形上定义（对照之下，未扭曲的超荷带旋量指标，只能定义在具有相应
%  自旋结构的流形上）。同一机制也解释了 $\chi_{\mu\nu}$（$\mathbf 3$ 分量）与
%  $\eta$（$\mathbf 1$ 分量）的来源。$Q$ 的 $\mathfrak u(1)_R$ 荷为 $+1$；$Q^2=0$ 来自
%  $\{\bar Q,\bar Q\}=0$——它是 (4.7) 第一式经共轭（或由 (4.7) 第二式的同族关系）得到的。]

现在，有了拓扑扭曲，我们就能在上文第 3 节的 CohTFT 范式中认出各个场：
\begin{enumerate}[label=\arabic*., leftmargin=2.2em]
  \item $A_{\mu}$ 仍是联络。但从场 $\psi^{A}_{\alpha}$ 得到一个奇 $1$-形式 $\psi_{\mu}$，
    且 $QA_{\mu}=\psi_{\mu}$。这些正对应第 3 节的场对 $(x^{i},\psi^{i})$。

  \item 由 $\bar{\psi}^{A}_{\dot\alpha}$，$\mathbf{2}\otimes\mathbf{2}$ 的对称积给出一个
    鬼数为 $-1$ 的奇自对偶形式 $\chi_{\mu\nu}$。这就是前面出现过的同一个 $\chi_{a}$。

  \item 由 $D_{\mu\nu}$ 得到一个鬼数为 $0$ 的偶自对偶形式。这“基本上”就是前面出现过
    的场 $H_{a}$。我们说“基本上”，是因为在物理理论中，经高斯消去后标准的辅助场 $D$
    被置为零（在没有物质超多重态时），而在一般的 CohTFT 形式中 $H$ 正比于方程。

  \item 由 $\bar{\psi}^{A}_{\dot\alpha}$ 还可以从 $\mathbf{2}\otimes\mathbf{2}$ 的反称积
    得到一个鬼数为 $-1$ 的奇零形式 $\eta$。它与 $\bar\phi$ 配对，给出构造
    $\Psi_{\mathrm{proj}}$ 所需的“投影多重态”。

  \item 剩下的是 $\phi$，它扮演 $\mathcal{G}$ 的李代数的对称代数的二次生成元的角色，
    这正是等变上同调的嘉当模型中所用的。
\end{enumerate}

% ---------------------------------------------------------------------
%  §4.2.2（续，原 PDF 第 25 页）
% ---------------------------------------------------------------------
\begin{enumerate}[label=\arabic*., leftmargin=2.2em, start=6]
  \item 也可以把 BRST 规范固定（包括通常的鬼场 $b,c$）考虑进来。这对应于为投影规范
    费米子选用另一个模型：用等变上同调的 Weil 模型而非嘉当模型。细节见 \cite{ref4}。
\end{enumerate}

% [Tutor Note: 这一条把 §3 的两种等变上同调模型落实到了具体的规范固定方案上：嘉当模型
%  （只用 $\phi$ 作等变参数）给出不含鬼场的 $\Psi_{\mathrm{proj}}$；Weil 模型（含
%  $b,c$）则给出标准 BRST 规范固定项。两模型在等变上同调层面等价（有理同伦论中的
%  “等变上同调模型定理”），故物理结果只差一个 $Q$-恰当项。这正是 (4.17) 中的
%  $\Psi$ 可以与 $\Psi_{\mathrm{loc}}+\Psi_{\mathrm{proj}}$ 不同而不影响结论的原因。]

现在，以这种方式转写各场之后，我们发现：把 CohTFT 形式（取 $s(A)=F^{+}$）给出的作用量
与物理作用量相比较，得到\footnote{\cite{ref4} 的 §8.11.3 指出 $\Psi$ 相对
$\Psi_{\mathrm{loc}}+\Psi_{\mathrm{proj}}$ 还有一个额外的移位；这相应于 $1$-形式
$V^{\dagger}$ 的一个移位，以计入嘉当复形上的规范作用。}\footnote{〔原稿边注〕这条
脚注现在需要修改！}
\begin{equation}
  S_{\mathrm{phys}}=Q(\Psi)+2\pi\ii\int\tau_{0}\,\tr(F\wedge F),
  \tag{4.17}
\end{equation}
其中
\begin{equation}
  \tau_{0}=\frac{\theta_{0}}{2\pi}+\frac{4\pi\ii}{g_{0}^{2}},
  \tag{4.18}
\end{equation}
而迹的归一化使得 $\int_{X}\tr(F\wedge F)$ 在任意 $X$ 上都是整数，且对 $X$ 上适当的
主 $G$-丛 $P$ 可取遍所有整数值。注意在推导 (4.17) 时用到了恒等式
\begin{equation}
  \int_{X}\tr(F^{\pm})^{2}=2\int_{X}\bigl(\tr F\ast F\pm\tr FF\bigr),
  \tag{4.19}
\end{equation}
这一点学生应当自行验证！

% [Tutor Note: (4.17) 与 (4.12) 的对应是本节最易读错的地方。(4.12) 中分别出现
%  $-\frac{1}{g_0^2}\int\tr(F\wedge\ast F)$ 与 $\frac{\theta_0}{8\pi^2}\int\tr(F\wedge F)$
%  两项，(4.17) 则把它们以及其余非拓扑项一并收进 $Q(\Psi)$，只留下度量无关的拓扑项
%  $2\pi\ii\tau_0\int\tr(F\wedge F)$，而把耦合常数组合成复耦合 $\tau_0$（(4.18)）。
%  这里的迹必须与 (4.12) 中同一个 $\tr$：只有 $\int\tr(F\wedge F)\in\mathbb{Z}$ 的整性
%  成立，$\tau_0\to\tau_0+1$（即 $\theta_0\to\theta_0+2\pi$）才与瞬子数平移
%  $k\to k+1$ 等价，$\theta$ 角的 $2\pi$ 周期性才有意义。若把 (4.12) 的 $\tr$ 换成
%  基本表示的 $\tr_{\mathrm{fund}}$，则 (4.18) 的系数必须改，否则 $\theta$ 周期会错。
%  另外请注意 (4.12) 的 $\frac{\theta_0}{8\pi^2}$ 与 (4.17) 中 $\tau_0$ 的
%  $\frac{\theta_0}{2\pi}$ 之间隔着一个 $2\pi\ii$ 的整体因子与 (4.19) 的恒等式，不宜
%  直接把两式的系数逐项对比。]

% [Tutor Note: (4.19) 的验证思路：按 (4.2) 的约定 $F^{+}=F+\ast F$（无 $1/2$），
%  $F^{-}=F-\ast F$。利用四维中 $2$-形式上的 $\ast^{2}=1$，有 $\ast F\wedge\ast F=F\wedge F$，
%  于是
%  $(F^{+})^{2}=(F+\ast F)^{2}=2(F\wedge F+F\wedge\ast F)$，
%  $(F^{-})^{2}=(F-\ast F)^{2}=2(F\wedge F-F\wedge\ast F)$。
%  取 $\tr$ 并积分即得 $\int\tr(F^{+})^{2}=2\int(\tr F\ast F+\tr FF)$，与 (4.19) 的
%  上号逐字相符。下号则需注意：由上式得
%  $\int\tr(F^{-})^{2}=2\int(\tr FF-\tr F\ast F)=-2\int(\tr F\ast F-\tr FF)$，与 (4.19)
%  下号相差一个整体符号（等价地，若把 $F^{-}$ 定义为 $\ast F-F$ 则无此问题）。因此
%  (4.19) 的 $\pm$ 应理解为“与 $F^{\pm}$ 同号”的简写；由于 $s(A)=F^{+}$，推导 (4.17)
%  时用到的主要是上号。这个整体符号不影响 $Q$-恰当项，但学生自己动手验证时容易在此处
%  卡住。]

现在我们可以对“当 $\omega_R=\omega^+$ 时 $Z(\omega^-,\omega^+,\omega_R)$ 与度量无关”
这一论断给出形式证明：在该轨迹（locus）上，QFT 具有 $Q$-对称性，而且除 (4.17) 中那个
与度量无关的项之外，作用量是 $Q$-恰当的。因此度量 $g_{\mu\nu}$ 的改变只导致作用量的
$Q$-恰当扰动，不应改变路径积分 $Z$，也不应改变 $Q$-闭算子的关联函数。

换一种说法：在物理上，对度量改变的响应由能量--动量张量刻画，但在目前情形
\begin{equation}
  T_{\mu\nu}:=\frac{1}{\sqrt{\det g_{\mu\nu}}}
  \frac{\delta}{\delta g^{\mu\nu}}S_{\mathrm{phys}}
  =\{Q,\Lambda_{\mu\nu}\}
  \tag{4.20}
\end{equation}
是 $Q$-恰当的。于是 $T_{\mu\nu}$ 与 $Q$-闭算子的关联函数应当为零。当我们考虑
$b_{2}^{+}(X)=1$ 的流形上算子关联函数的穿墙现象时，将需要重新审视这一论证。
\footnote{〔原稿边注〕从第 4 节开始已经过去 34 分钟。}

% ---------------------------------------------------------------------
%  §4.3
% ---------------------------------------------------------------------
\subsection{瞬子模空间简述}
\label{sec:4-3}

若接受一般局域化公式 (3.54) 也适用于无穷维路径积分，那么它断言：规范群为 $G$ 的矢量
多重态的 SYM 路径积分将局域化到 $s(A)$ 的零点模空间上的积分，即反自对偶（ASD）联络
（又称瞬子）的模空间：
\begin{equation}
  \mathcal{M}(P,g)=\{A\in\mathcal{A}\,|\,F^{+}=0\}/\mathcal{G}.
  \tag{4.21}
\end{equation}
下面我们就来稍作讨论。

\subsubsection{形变复形与虚拟维数}
\label{sec:4-3-1}

先考虑模空间的切空间。全体联络构成的空间在 $A$ 处的切空间为
$T_{A}\mathcal{A}\cong\Omega^{1}(\mathrm{ad}P)$；也就是说，过 $A$ 的任一条联络路径在
$t$ 的一阶上都形如 $A+t\alpha$，其中 $\alpha$ 是取值于 $P$ 的伴随丛的整体定义
$1$-形式。简单的计算表明
\begin{equation}
  F^{+}(A+t\alpha)=F^{+}(A)+t\nabla^{+}_{A}\alpha+O(t^{2}).
  \tag{4.22}
\end{equation}
这就定义了算子 $\nabla^{+}_{A}:\Omega^{1}(\mathrm{ad}P)\to\Omega^{2,+}(\mathrm{ad}P)$。
现在，若考虑模空间内的一条联络路径，则切向量 $\alpha$ 必须属于该算子的核；另一方面，
在除以规范群之前，由规范变换得到的、$F^{+}(A)=0$ 的解空间在 $\mathcal{A}$ 中的切向量
是平凡的。上述考虑促使我们定义瞬子方程的一阶形变复形：
\begin{equation}
  0\to\Omega^{0}(\mathrm{ad}P)
  \xrightarrow{\ \nabla_{A}\ }
  \Omega^{1}(\mathrm{ad}P)
  \xrightarrow{\ \nabla^{+}_{A}\ }
  \Omega^{2}(\mathrm{ad}P)^{+}.
  \tag{4.23}
\end{equation}

% [Tutor Note: (4.23) 只有在 $F^{+}(A)=0$ 的解处才是复形：一般地
%  $\nabla^{+}_{A}\nabla_{A}\alpha=[F^{+}(A),\alpha]$，故复合为零恰好等价于瞬子方程。
%  这也解释了为何只有该复形的上同调才给出模空间的切空间：$H^{0}$ 衡量可约性
%  （规范等价类维数），$H^{1}$ 是形变，$H^{2}$ 是阻碍；虚拟维数
%  $\vdim=H^{1}-H^{0}-H^{2}$。]

至少形式上，该复形的上同调应当给出瞬子模空间的切空间。事实上，正是出于这个原因，
阿蒂亚、希钦与辛格在 \cite{ref1} 中研究了这一复形。他们证明它是椭圆的，并且
如果把复形卷起来，得到的恰好是手征狄拉克算子
\begin{equation}
  D^{+}:\Gamma(S^{-}\otimes\mathcal{V})\to\Gamma(S^{+}\otimes\mathcal{V})
  \tag{4.24}
\end{equation}
耦合到复向量丛
\begin{equation}
  \mathcal{V}=\mathrm{ad}P\otimes S^{+}
  \tag{4.25}
\end{equation}
上的情形。这表明 ASD 瞬子的模空间 $\mathcal{M}(P,g)$ 具有虚拟维数
\begin{equation}
  \begin{aligned}
    \vdim\mathcal{M}(P,g)
      &=4h^{\vee}k-\dim G\,\bigl(b_{2}^{+}-b_{1}+1\bigr)\\
      &=4h^{\vee}k-\dim G\,\frac{\chi+\sigma}{2},
  \end{aligned}
  \tag{4.26}
\end{equation}
其中 $h^{\vee}$ 是 $\mathfrak{g}$ 的对偶考克斯特数，而
\begin{equation}
  k:=\frac{1}{16\pi^{2}h^{\vee}}\int_{X}\Tr_{\mathfrak{g}}F^{2}
  =p_{1}(\mathrm{ad}P)/4h^{\vee}
  \tag{4.27}
\end{equation}
是整数（当 $\mathfrak{g}$ 为单李代数时第二个等号成立）。我们选取 $k$ 的归一化，使得当
$X=S^{4}$ 时对适当的丛 $P$ 可取遍所有整数值。符号取得使反自对偶瞬子对应 $k>0$。
\footnote{〔原稿边注〕需要核对第二个等号处的符号。}特别地，对 $SU(2)$ 有
\begin{equation}
  \vdim\mathcal{M}(P,g)=8k-3\bigl(b_{2}^{+}-b_{1}+1\bigr),
  \tag{4.28}
\end{equation}
其中
\begin{equation}
  k=-\frac{1}{8\pi^{2}}\int_{X}\Tr_{2}F^{2}.
  \tag{4.29}
\end{equation}

% [Tutor Note: (4.26) 的两个等号用了 $\chi=2-2b_1+b_2$、$\sigma=b_2^+-b_2^-$、
%  $b_2=b_2^++b_2^-$，故 $(b_2^+-b_1+1)=(\chi+\sigma)/2$。常数项
%  $\dim G\,(b_2^+-b_1+1)$ 的来源是：$H^{0}$ 与 $H^{2}$ 的维数分别为 $\dim G$ 与
%  $\dim G\,b_2^+$，指标 $=$ (1 形式零模数) $-$ (0 形式) $-$ (2 形式)
%  $=8h^\vee k-\dim G-\dim G\,b_2^+$，再加上规范固定与 $H^{1}$ 中调和形式的
%  $\dim G\,b_1$ 修正，合起来正好是上式。]

% [Tutor Note: (4.27) 与 (4.29) 的符号约定并不一致，这恰是作者边注“需要核对第二个等号
%  处的符号”所指之处。对 $SU(2)$，$h^\vee=2$，李代数迹与基本表示迹满足
%  $\Tr_{\mathfrak g}F^2=2h^{\vee}\Tr_{2}F^2=4\Tr_{2}F^2$；代入 (4.27) 得
%  $k=+\frac{1}{8\pi^2}\int\Tr_2F^2$，而 (4.29) 给的是
%  $k=-\frac{1}{8\pi^2}\int\Tr_2F^2$，两者相差整体符号。正文“符号取得使 ASD 瞬子
%  $k>0$”这一要求只与 (4.29) 相容：ASD 时 $F=-\ast F$，故
%  $\Tr_2F^2=-\Tr_2(F\ast F)$ 为负，于是 (4.29) 给出 $k>0$。使用 (4.27) 时须记住
%  这里的 $\Tr_{\mathfrak g}$ 归一化与下文 §4.4 使用的 $\Tr_2$ 相差一个符号。]

\subsubsection{虚拟维数就是鬼数反常}
\label{sec:4-3-2}

现在把这个讨论与 CohTFT 作比较是很有意思的。复形 (4.23) 是一般讨论中序列 (3.52) 的
特例；特别地，那里的算子 $\mathbb{F}$ 确实是弗雷德霍姆算子。

此外，矢量多重态中的费米子带 $u(1)_R$ 荷 $\pm1$。因此在把 $S_{R}\cong S^{+}$ 认同起来的
拓扑扭曲之后，那个其指标计算 $u(1)_R$ 对称性反常的狄拉克算子，恰好就是用来计算 AHS
复形 (4.24) 指标的同一个算子。现在回忆：$u(1)_R$ 荷就是鬼数——形式上等同于
$\mathcal{M}(P,g)$ 上微分形式的次数。于是：

\begin{center}
\itshape 指标 $\Ind(\mathbb{F})$ 与 $u(1)_R$ 对称性的反常一致，也等于鬼数反常。
\end{center}

% [Tutor Note: 这一等同是全书的关键等式之一，值得停下来想清楚其结构。AHS 复形 (4.23) 的
%  指标（按复形取 Euler 特征，即 $\sum_i(-1)^i\dim H^i$）由 (4.24) 的“卷起”同构等于手征
%  狄拉克算子 $D^{+}:\Gamma(S^-\otimes\mathcal V)\to\Gamma(S^+\otimes\mathcal V)$ 的指标，
%  故它由阿蒂亚--辛格指标定理给出：$\Ind(D^{+})=\int_X\hat A(X)\,
%  \mathrm{ch}(\mathcal V)$，而 $\mathcal V=\mathrm{ad}P\otimes S^{+}$ 使
%  $\mathrm{ch}(\mathcal V)=\mathrm{ch}(\mathrm{ad}P)\,\mathrm{ch}(S^{+})$。把
%  $\hat A\,\mathrm{ch}(S^{+})$ 合并即得 (4.26) 的表达式（整体符号取决于复形指标的取向
%  约定，须与 (4.27)(4.29) 的符号约定一并核对）。另一方面，$u(1)_R$ 是
%  $\mathcal{N}=2$ 的 $R$-对称，其流的三角图反常只由带荷 $\pm1$ 的费米子
%  （$\psi,\bar\psi$）贡献；扭曲后它们正是 $S^{\pm}\otimes\mathrm{ad}P$ 的截面，故反常
%  与上面这个指标是同一个数。鬼数则是把 $u(1)_R$ 荷赋给场后、$Q$ 使其降一阶所得的分次；
%  路径积分不湮灭要求插入算子的总鬼数恰好抵消反常，这正是 (4.45)(4.58) 的来源。]

\subsubsection{瞬子模空间的一些基本数学性质}
\label{sec:4-3-3}

\begin{enumerate}[label=\arabic*., leftmargin=2.2em]
  \item 模空间 $\mathcal{M}(P,g)$ 一般是奇异的、非紧的。当 (4.23) 的 $H^{0}$ 非零且
    联络可约、或 $H^{2}$ 非零而瞬子的形变存在成为真正形变的阻碍时，就会出现奇点。
    然而，一般度量定理 \cite{ref15} 断言：当 $k>0$ 且 $G=SU(2)$ 或 $G=SO(3)$ 时，只要
    $b_{2}^{+}(X)>0$，对一般的度量模空间就是光滑的。因此在上文 CohTFT 的一般局域化
    公式 (3.54) 中，将有 $\operatorname{cok}(\mathbb{F})=0$，无需引入欧拉类。在 \cite{ref22} 中，
    克龙海默把一般度量定理的某些方面推广到了 $PSU(N)$ 丛的情形。
    \footnote{〔原稿边注〕应当再多说一些。}

  \item 在光滑点处，
    \begin{equation}
      \vdim\mathcal{M}(P,g)=\dim\mathcal{M}(P,g)
      \tag{4.30}
    \end{equation}
    就是模空间 $\mathcal{M}(P,g)$ 的真实维数。

  \item 模空间 $\mathcal{M}(P,g)$ 还是非紧的，因为瞬子可以退化为点瞬子，此时
    $\Tr F\wedge F$ 是一个光滑测度加上若干个狄拉克测度之和。
\end{enumerate}

\subsubsection{瞬子模空间作为 Q 不动点轨迹}
\label{sec:4-3-4}

由于下文要讨论纯矢量多重态理论的推广，把它看作一个 $Q$ 不动点方程是值得的。我们有：
\begin{equation}
  Q\chi_{\mu\nu}=\ii\bigl(F^{+}_{\mu\nu}-D^{+}_{\mu\nu}\bigr).
  \tag{4.31}
\end{equation}
通过高斯积分消去辅助场给出 $D^{+}_{\mu\nu}=0$。当我们在下文引入超多重态后，方程
$D^{+}=0$ 将发生非常重要的改变。

有意思的是，理论中其余费米子的 $Q$ 不动点方程标出了麻烦应当在何处出现：令
$Q\psi_{\mu}=0$ 给出
\begin{equation}
  D_{A}\phi=0,
  \tag{4.32}
\end{equation}
而 $Q\eta=0$ 给出
\begin{equation}
  [\bar\phi,\phi]=0.
  \tag{4.33}
\end{equation}
只有当联络不可约时，方程 (4.32) 才仅有解 $\phi=0$。在可约联络处，我们有诸多理由预期
出现麻烦。方程 (4.33) 在低能有效理论（LEET）的物理分析中起重要作用，下文将予以说明。

% ---------------------------------------------------------------------
%  §4.4
% ---------------------------------------------------------------------
\subsection{可观测（量）}
\label{sec:4-4}

为得到有意思的路径积分，我们需要插入 $Q$-闭且规范不变的观测量。回忆在嘉当模型中
\begin{equation}
  Q\phi=0 .
  \tag{4.34}
\end{equation}
现在 $\phi$ 取值于伴随表示，故不是规范不变的；然而给定 $\mathfrak{g}$ 上任一不变多项式
$\mathcal{P}$，就可以构造出 $Q$-闭且规范不变的可观测量 $\mathcal{P}(\phi)$。我们称
这些为 0-可观测量，因为它们是定义在\textbf{点}上的局部算子。它们构成一个环，由
$\mathfrak{g}$ 的卡西米尔生成。例如对 $SU(N)$，有生成观测量
\begin{equation}
  \mathcal{O}^{(0)}_{s}(\wp):=\Tr\phi^{s}(\wp),\qquad s=2,\dots,N .
  \tag{4.35}
\end{equation}

给定任一零维可观测量 $\mathcal{O}^{(0)}$，可以利用下降（descent）形式构造一族非局域
可观测量。新的可观测量可以如下典范地构造：在拓扑扭曲下
$Q_{\alpha\dot\alpha}\to K_{\mu}$，且
\begin{equation}
  \{Q,K_{\mu}\}=\partial_{\mu}.
  \tag{4.36}
\end{equation}
于是若令 $K:=\dd x^{\mu}K_{\mu}$ 为取值 $1$-形式的超对称算子，便可定义
$\mathcal{O}^{(1)}:=K\mathcal{O}^{(0)}$ 而得到取值 $1$-形式的算子。现在，对任一
$1$-链 $\gamma$ 有
\begin{equation}
  Q\int_{\gamma}\mathcal{O}^{(1)}
  =\mathcal{O}^{(0)}\big|_{\partial\gamma}.
  \tag{4.37}
\end{equation}
这意味着：
\begin{enumerate}[label=\arabic*., leftmargin=2.2em]
  \item 点 $\wp$ 在 $\mathcal{O}^{(0)}(\wp)$ 中位置的改变是 $Q$-恰当的：若
    $\partial\gamma=\wp_{1}-\wp_{2}$，则
    $\mathcal{O}^{(0)}(\wp_{1})=\mathcal{O}^{(0)}(\wp_{2})
    +Q\int_{\gamma}\mathcal{O}^{(1)}$。
  \item 若 $\gamma$ 是闭链，则 $\mathcal{O}(\gamma):=\int_{\gamma}\mathcal{O}^{(1)}$
    是 BRST 闭的。
\end{enumerate}
类似地，$\mathcal{O}^{(j)}:=K^{j}\mathcal{O}^{(0)}$ 定义 $X$ 上的 $j$-形式；若
$\Sigma_{j}$ 是闭的 $j$-链，则
\begin{equation}
  \mathcal{O}(\Sigma_{j}):=\int_{\Sigma_{j}}\mathcal{O}^{(j)}
  \tag{4.38}
\end{equation}
是 $Q$-闭、规范不变且只依赖于 $\Sigma_{j}$ 的同调类的可观测量。我们称它们为
“$j$-可观测量”。

% [Tutor Note: (4.37) 的证明只用 (4.36) 与斯托克斯定理；注意 $K_\mu$ 是\textbf{超对称}
%  算子（把超荷按 $\sigma^\mu$ 的指标结构重组而成），它把鬼数提高 $1$ 并把形式次数提高
%  $1$。因此 (4.38) 中 $\mathcal{O}(\Sigma_j)$ 的总鬼数等于 $\mathcal{O}^{(0)}$ 的鬼数
%  加上 $j$，而非仅由零维算子决定——这正是 (4.45) 中 $4\ell+2r$ 的来历
%  （$\mathcal{O}$ 的鬼数 $4$、$\mathcal{O}(\Sigma)$ 的鬼数 $2$）。]

在下文的计算中我们主要关注秩一的群 $SU(2)$ 与 $SO(3)$。此时所有 0-可观测量都由那个
鬼数 $=4$ 的可观测量生成：
\begin{equation}
  \mathcal{O}^{(0)}_{2}(\wp):=\frac{1}{8\pi^{2}}\Tr_{2}\phi^{2}(\wp).
  \tag{4.39}
\end{equation}

至此我们采用如下策略：为使公式易读，某些方程中我们将省略实系数；它们通常是 $2$ 与
$\pi$ 的（分数次）幂。这样做时我们用符号 $\sim$；而当我写“$=$”时，我确实是指“等于”。
带正确系数的完整表达式可在 \cite{ref38} 中找到。

在秩一拓扑扭曲 SYM 中特别重要的是两个可观测量，它们算出来是
\begin{equation}
  \mathcal{O}^{(2)}_{2}(\Sigma)\sim\int_{\Sigma}\Tr\bigl(\phi F+\psi\wedge\psi\bigr).
  \tag{4.40}
\end{equation}

为简化记号，此后我们写：
\begin{equation}
  \mathcal{O}^{(0)}_{2}(\wp)\to\mathcal{O}.
  \tag{4.41}
\end{equation}
注意由于对点 $\wp$ 的依赖是 $Q$-恰当的，可以从记号中去掉位置 $\wp$。类似地我们直接写
\begin{equation}
  \mathcal{O}^{(2)}_{2}(\Sigma)\to\mathcal{O}(\Sigma).
  \tag{4.42}
\end{equation}

现在考虑这些算子在拓扑扭曲理论（我们把它视为 $\omega_{R}=\omega^{+}$ 的物理理论）中的
关联函数。于是我们的记号是：对场中任一表达式 $\mathcal{F}$（通常是算子的乘积），记
\begin{equation}
  \Bigl\langle\prod_{i}\mathcal{F}_{i}\Bigr\rangle_{\mathcal{T}}
  :=\int[\dd A\,\dd\psi\cdots]\,\ee^{S_{\mathcal{T}}}\,\mathcal{F}
  \tag{4.43}
\end{equation}
表示在 $\omega_{R}=\omega^{+}$ 下计算的 $\mathcal{N}=2$ 场论 $\mathcal{T}$ 的路径积分。

下面把“威滕多项式”定义为同调 $H_{0}(X)\oplus H_{2}(X)$ 上带实系数的多项式函数：
\begin{equation}
  P_{W}(\wp^{\ell}\Sigma^{r})
  :=\bigl\langle\mathcal{O}^{\ell}\mathcal{O}(\Sigma)^{r}\bigr\rangle_{\mathcal{T}}.
  \tag{4.44}
\end{equation}
我们注意到：
\begin{enumerate}[label=\arabic*., leftmargin=2.2em]
  \item 它们是同调上形式地与度量无关、从而只依赖于光滑结构的复值多项式。
  \item 由于鬼数反常，$P_{W}(\wp^{\ell}\Sigma^{r})$ 只有在
    \begin{equation}
      4\ell+2r=\vdim\mathcal{M}(P,g)
      \tag{4.45}
    \end{equation}
    时才可能非零。
\end{enumerate}

\textbf{评注：}

\begin{enumerate}[label=\arabic*., leftmargin=2.2em]
  \item 特别注意 $\eta$ 与 $\chi$ 都不出现在可观测量中。这在下文 $u$-平面积分的计算中
    将起非常重要的作用。

  \item 我们强调：即使做了拓扑扭曲（即令 $\omega_{R}=\omega^{+}$），$X$ 上几乎所有算子
    的关联函数仍将依赖度量（即依赖自旋联络 $\omega^{\pm}$）。例如若考虑
    $\Tr F_{\mu\nu}F^{\mu\nu}$ 或 $\Tr(D_{\mu}F^{\mu\nu})^{2}$ 之类的关联函数，就会有
    度量依赖。物理学家所说的“理论的一个拓扑扇区”，指的正是只考虑上述这类 $Q$-闭
    可观测量。
\end{enumerate}

% [Tutor Note: 第 1 条评注是 §5 中 $u$-平面积分能被算出且只依赖拓扑数据的根本原因：
%  $\eta,\chi$ 不出现在插入算子中，故在局域化到 $u$-平面上的常数 $Q$ 不动点后，只剩下
%  $\phi$（及其下降形式）的幂，而 $\phi$ 在同调层面被识别为模空间上的 $\Theta$-曲率
%  （见 (4.53) 及其后的说明）。等价地说，$\mathcal{O}$ 与 $\mathcal{O}(\Sigma)$ 生成
%  的正是唐纳森的 $\mu$ 映射的像。]

% ---------------------------------------------------------------------
%  §4.5
% ---------------------------------------------------------------------
\subsection{唐纳森多项式}
\label{sec:4-5}

先简要回顾唐纳森多项式。我们从唐纳森的 $\mu$ 映射出发：
\begin{equation}
  H_{*}(X)\to H^{*}(\mathcal{M}),
  \tag{4.46}
\end{equation}
即从 $X$ 的同调到 $\mathcal{M}$ 的上同调的映射。下面简述唐纳森的构造。

回顾唐纳森的表述：考虑主 $G$-丛及其分类映射 $f$：
\begin{equation}
  \begin{array}{c}
    P\times\mathcal{A}/\mathcal{G}\\[3pt]
    \big\downarrow\\[3pt]
    (P\times\mathcal{A})/(\mathcal{G}\times\mathcal{G})
    \ \xrightarrow{\ f\ }\ BG
  \end{array}
  \tag{4.47}
\end{equation}
若在 $\mathfrak{g}$ 上选取任一 $d$ 次不变多项式 $\mathcal{P}$，从而得到类
$\varpi\in H^{2d}(BG)$，则可定义斜积
\begin{equation}
  \mu_{D}(\Sigma_{j}):=\int_{\Sigma_{j}}f^{*}(\varpi)
  \in H^{2d-j}(\mathcal{A}/\mathcal{G}).
  \tag{4.48}
\end{equation}
这是一个上同类，可限制到 $\mathcal{M}\subset\mathcal{A}/\mathcal{G}$ 上，且只依赖
$\Sigma_{j}$ 的同调类。

对 $SU(2)$，取 $\varpi$ 为 $H^{4}(BSU(2);\mathbb{Z})\cong\mathbb{Z}$ 的生成元，
由此定义形式
\begin{equation}
  \wp\to\mu_{D}(\wp)\in H^{4}(\mathcal{M}),
  \tag{4.49}
\end{equation}
\begin{equation}
  \Sigma\to\mu_{D}(\Sigma)\in H^{2}(\mathcal{M}).
  \tag{4.50}
\end{equation}

现在回忆：唐纳森在 $H_{0}(X)\oplus H_{2}(X)$ 上定义他的多项式
\footnote{〔原稿边注〕表示曲面时大概用 $\Sigma$ 比用 $S$ 更好。}，办法是说明可以选取上同
调类的紧支集代表，并把单项式 $\wp^{\ell}\Sigma^{r}$ 上的值定义为
\begin{equation}
  P_{D}(\wp^{\ell}\Sigma^{r})
  :=\int_{\mathcal{M}}\mu_{D}(\wp)^{\ell}\,\mu_{D}(\Sigma)^{r}.
  \tag{4.51}
\end{equation}
也就是说，这些系数由模空间上的相交数给出。由严格的数学分析我们知道：

\begin{enumerate}[label=\arabic*., leftmargin=2.2em]
  \item $P_{D}(\wp^{\ell}\Sigma^{r})$ 是有理数，且与度量无关；例外是
    $b_{2}^{+}(X)\le1$ 的 $X$。当 $b_{2}^{+}(X)>1$ 时，对一般度量模空间
    $\mathcal{M}(P,g)$ 是光滑的，这就是 Freed--Uhlenbeck 的一般度量定理
    \cite{ref15}。因此在这些情形下它们给出 $X$ 的光滑不变量。当 $b_{2}^{+}(X)=1$ 时，
    这些多项式在度量空间中只是分段常数的，并会在实余维一的墙上发生跳跃。
    \footnote{当 $Q_{X}$ 正定时，可以选取定向使 $b_{2}^{+}(X)=0$。此时模空间不必光滑；
    特别地，取 $k=1$ 时它是 $X$ 与一组奇点之间的 $5$ 维配边，这些奇点对应于可约联络，
    在奇点附近模空间局部地是 $\CC\PP^{2}$ 或 $\overline{\CC\PP^{2}}$
    上的锥。唐纳森证明“$Q_{X}$ 正定时相交形式可在整数上对角化”用的正是这一点。}
\end{enumerate}

\textbf{评注：}依照博利厄与辛格的讨论，可以给出把 $\mathcal{O}$ 与
$\mathcal{O}(\Sigma)$ 联系到 $\mu_{D}(\wp)$ 与 $\mu_{D}(\Sigma)$ 的精确论证。它用到
$\mathcal{A}$ 的 $\mathcal{G}$-等变上同调的一个模型，以及
\begin{equation}
  P\times\mathcal{A}\to X\times\mathcal{A}/\mathcal{G}
  \tag{4.52}
\end{equation}
上的“泛联络”。细节见 \cite{ref4} §8.8。博利厄与辛格论证道，可以把
\begin{equation}
  \mathcal{F}=F+\psi+\phi
  \tag{4.53}
\end{equation}
认作泛联络的曲率，并按沿 $X$ 与沿 $\mathcal{A}/\mathcal{G}$ 的形式次数分解。于是
\begin{equation}
  \Tr\mathcal{F}^{2}
  \tag{4.54}
\end{equation}
上的斜积就把唐纳森的映射 $H_{*}(X)\to H^{*}(\mathcal{M}(P,g))$ 与下降形式给出的、
从 $H_{*}(X)$ 到 $Q$-上同调 $H^{*}_{Q}$ 的映射联系了起来。

当然，直接验证下述事实并不困难：在计算 $\phi$ 的量子关联函数时，可以把它认同为泛联络
曲率沿 $\mathcal{A}/\mathcal{G}$ 方向的分量。

在上文局域化公式 (3.54) 的最终表述中，我们有一个从规范不变的等变闭类
$\hat{\mathcal{O}}$ 到 $\mathcal{M}=Z(s)/\mathcal{G}$ 上微分形式
$\omega_{\hat{\mathcal{O}}}$ 的映射。把这一切合在一起，对 $\mathfrak{g}$ 上每个 $d$ 次
不变多项式，我们得到交换图
\begin{equation}
  \begin{array}{ccc}
    H_{*}(X) & \xrightarrow{\ \mathrm{descent}\ } & H^{2d-*}_{Q}\\[6pt]
     & \searrow_{\ \mu} & \big\downarrow{\scriptstyle
       \hat{\mathcal{O}}\mapsto\omega_{\hat{\mathcal{O}}}}\\[10pt]
     & & H^{2d-*}(\mathcal{M})
  \end{array}
  \tag{4.55}
\end{equation}

于是我们自然被引向威滕 1988 年论文的主要断言：

\begin{center}
\itshape 在上述对应之下，$Q$-闭算子的物理关联函数与唐纳森多项式重合——相差一个整体
常数。这正是威滕 1988 年论文的目标：0-可观测量与 2-可观测量分别精确对应于唐纳森的
形式 $\omega_{D}(\wp)$ 与 $\omega_{D}(\Sigma)$。
\end{center}

% [Tutor Note: (4.55) 的交换图是全节的形式总结：左边是 $X$ 的同调，右下是模空间的上同
%  调，两条路径分别是“下降 $+$ $Q$-上同调到微分形式的同构”与“$\mu_D$ 斜积”。图能交换
%  的实质是：$F+\psi+\phi$ 在等变意义下就是一个曲率，而 $\Tr\mathcal F^2$ 的斜积既是
%  唐纳森的 $\mu$ 映射、又是下降形式。注意 (4.54) 只用 $\Tr\mathcal F^2$（二次
%  卡西米尔）就足以生成 $SU(2)$ 的全部不变量，因为更高次卡西米尔对 $SU(2)$ 都是它的
%  多项式。]

% ---------------------------------------------------------------------
%  §4.6
% ---------------------------------------------------------------------
\subsection{唐纳森--威滕配分函数}
\label{sec:4-6}

在物理上，把一族关联函数组装成单个生成函数是极其自然的，它对应于用所讨论的可观测量
扰动理论的作用量。我们称这个生成函数为\textbf{唐纳森--威滕配分函数}：
\begin{equation}
  Z_{DW}(p,\Sigma):=\bigl\langle\ee^{p\mathcal{O}+\mathcal{O}(\Sigma)}
  \bigr\rangle_{\mathcal{T}}.
  \tag{4.56}
\end{equation}
将其形式展开可得
\begin{equation}
  Z_{DW}(p,\Sigma)
  =\sum_{\ell,r\ge0}\frac{p^{\ell}}{\ell!\,r!}
  \bigl\langle\mathcal{O}^{\ell}\mathcal{O}(\Sigma)^{r}\bigr\rangle_{\mathcal{T}}.
  \tag{4.57}
\end{equation}
固定 $\ell,r$ 的每个关联函数都是对瞬子扇区的求和：于是要对主 $SU(2)$ 丛 $P$ 的同构类
求和。然而由 $u(1)_R$ 反常给出的选择定则可知，只有满足
\begin{equation}
  4\ell+2r=\vdim\mathcal{M}(P,g)
  \tag{4.58}
\end{equation}
的瞬子扇区才可能贡献。于是可以把威滕的主要断言强化如下：对由 $SU(2)$ SYM 理论给出的
$\mathcal{T}$，
\begin{equation}
  \begin{aligned}
    Z_{DW}(p,s)
    &:=\bigl\langle\ee^{p\mathcal{O}+\mathcal{O}(\Sigma)}\bigr\rangle_{\mathcal{T}}
    =\frac{1}{2}\Lambda^{-\frac{3}{4}(\chi+\sigma)}
    \sum_{\ell,r\ge0}\frac{p^{\ell}s^{r}}{\ell!\,r!}
    \Lambda^{2\ell+|r|}P_{D}(\wp^{\ell}\Sigma^{r}).
  \end{aligned}
  \tag{4.59}
\end{equation}

% [Tutor Note: (4.59) 中的 $\Lambda$ 是量子场论路径积分里隐含的标度参数，定义量子理论
%  必需它，不能简单置 $1$。理由有三：(i) 经典作用量无量纲，但重整化后 $g_0$ 变成随标度
%  跑的 $g(\mu)$（$\mathcal{N}=2$ 中 β 函数单圈精确，见 (4.76)），$\Lambda$ 正是
%  $\mathcal{N}=2$ SYM 的动力学标度 $\Lambda^{2h^\vee}=\mu^{2h^\vee}
%  \ee^{-8\pi^{2}/g^{2}(\mu)}$（$SU(2)$ 时 $h^\vee=2$，即 $\Lambda^{4}$），故 $\Lambda$
%  记录 UV 与 IR 的匹配；(ii) 瞬子展开的权重是 $\Lambda^{2h^\vee k}$ 型（$SU(2)$ 时即
%  $\Lambda^{4k}$），若置 $\Lambda=1$ 就丢掉了 $k$ 依赖的标度；(iii) 一致性检验：由
%  (4.26)(4.58)，$4\ell+2r=4h^\vee k-\dim G\,(b_2^+-b_1+1)$，故
%  $2\ell+|r|=\vdim/2=2h^\vee k-\frac{\dim G}{4}(\chi+\sigma)$，其中 $|r|$ 实为
%  $\sum r_\alpha$ 的简写（见 (4.61)）；把它与前置因子
%  $\Lambda^{-\frac34(\chi+\sigma)}=\Lambda^{-\frac{\dim G}{4}(\chi+\sigma)}$
%  （此处 $\dim G=3$）相乘，
%  总幂次恰为 $\Lambda^{2h^\vee k-\frac{\dim G}{2}(\chi+\sigma)}$，即“瞬子权重
%  $\times$ 流形依赖的指数因子”这一标准结构。把 $\Lambda$ 置 $1$ 会把这两重结构同时
%  抹平，使 (4.59) 退化为纯组合恒等式而看不出标度内容——这正是作者说“误导”的意思。]

此处须作几点说明：

\begin{enumerate}[label=\arabic*., leftmargin=2.2em]
  \item 我们引入了一个参数 $\Lambda$，它是量子场论路径积分中隐含的。为定义量子理论它
    是必需的，下文还将对它作更多评论。在物理进路讨论唐纳森理论的大多数文献中它被置为
    $1$，但这实际上是误导性的。

  \item (4.59) 中的路径积分是 $\mathcal{O}$ 与 $\mathcal{O}(\Sigma)$ 关联函数的生成
    泛函。如同 $SU(2)$ 规范理论中的所有路径积分一样，我们必须对带联络的丛求和；对丛的
    求和也称为对瞬子扇区求和。对给定的插入 $\mathcal{O}^{\ell}\mathcal{O}(\Sigma)^{r}$，
    丛求和中只有一项会贡献，即满足鬼数选择定则 (4.45) 的那一项。

  \item 数学上的讲究者会坚持认为：向量空间上的多项式取值于该空间对偶的对称代数。确实，
    我们应把 $p$ 看作同调类 $[\wp]\in H_{0}(X)$ 的对偶。按此精神，$\mathcal{O}(\Sigma)$
    应解释如下：取 $H_{2}(X;\mathbb{Z})$ 的一组基 $[\Sigma_{\alpha}]$ 及对偶基
    $s^{\alpha}$，则
    \begin{equation}
      \mathcal{O}(\Sigma)=s^{\alpha}\mathcal{O}(\Sigma_{\alpha}).
      \tag{4.60}
    \end{equation}
    于是生成函数 $Z_{DW}(p,\Sigma)$ 成为 $p$ 与 $s^{\alpha}$ 的形式幂级数。对 $r\ge0$
    的求和是对多重指标 $\vec r=(r_{1},\dots,r_{b_{2}})$ 求和，且
    $|r|=\sum r_{\alpha}$；因此我们的简写意为
    \begin{equation}
      \sum_{r\ge0}\Lambda^{|r|}\frac{s^{r}}{r!}\mathcal{O}(\Sigma)^{r}
      :=\sum_{r_{\alpha}\ge0}\prod_{\alpha}
      \frac{(\Lambda s_{\alpha})^{r_{\alpha}}}{r_{\alpha}!}
      \mathcal{O}(\Sigma_{\alpha})^{r_{\alpha}}.
      \tag{4.61}
    \end{equation}
    级数中任何一项的系数都是良定义的，但并无任何断言说整个和在 $H_{*}(X)$ 的同调上作为
    函数有意义，而不只是作为 $H_{*}(X)$ 上多项式的生成级数。我们常常对此颇为马虎，
    径自称 $Z_{DW}(p,\Sigma)$。

  \item 我们处理的是 $SU(2)$ 规范理论，但场处于伴随表示中，故对“扭曲的 $SU(2)$ 丛”
    （即不提升为 $SU(2)$ 丛的 $SO(3)$ 丛）定义路径积分也是有意义的。$X$ 上的主
    $SO(3)$ 丛有两个示性类：$\xi=w_{2}(P)\in H^{2}(X;\mathbb{Z}_{2})$（我们称之为
    't Hooft 磁通）与瞬子数 $k\in H^{4}(X;\mathbb{Z})\cong\mathbb{Z}$。在生成函数中
    我们对瞬子数 $k$ 求和（这对 $SU(2)$ 规范理论的路径积分是恰当的），但当 't Hooft
    磁通非零时，把路径积分推广到取值于 $\mathrm{ad}P$ 的场并无困难。相应的生成函数
    记作
    \begin{equation}
      Z_{DW}(p,\Sigma)\to Z^{\xi}_{DW}(p,\Sigma).
      \tag{4.62}
    \end{equation}
    我们强调：我们做的是 $SU(2)$ 规范理论——对瞬子数求和，但不对 $\xi$ 求和。在
    $SO(3)$ 规范理论中我们还必须对 $\xi$ 求和。

  \item 该生成函数与克龙海默和姆罗夫卡在 \cite{ref21} 中引入的生成函数密切相关：
    \begin{equation}
      Z^{\xi}_{KM}(p,s):=\sum_{\ell,r\ge0}\frac{p^{\ell}s^{r}}{\ell!\,r!}
      P_{D}(\wp^{\ell}\Sigma^{r}).
      \tag{4.63}
    \end{equation}
    相对的整体因子 $\frac{1}{2}$ 来源于物理学家要除以 $SU(2)$ 中心的阶，而该中心并不
    有效地作用于场。

  \item 作为一个简单例子，对 $X=K3$（通常的 K3 曲面），可取 $\xi=0$
    \footnote{〔原稿边注〕合适的参考文献？Kotschick--摩根。}，此时
    \begin{equation}
      Z_{KM}(p,s)=\sinh\Bigl(\frac{1}{2}s^{2}+2p\Bigr).
      \tag{4.64}
    \end{equation}

  \item 当然，当 $X$ 非单连通时，可以把 $1$-可观测量也包括进来。
    \footnote{〔原稿边注〕本节到此结束：从 §4.3 起已过去 52 分钟。}
\end{enumerate}

% ---------------------------------------------------------------------
%  §4.7
% ---------------------------------------------------------------------
\subsection{推广：拉氏量型 $\mathcal{N}=2$ 理论}
\label{sec:4-7}

规范群为 $SU(2)$ 的 $\mathcal{N}=2$ SYM 理论只是一大类 $\mathcal{N}=2$ QFT 中的一个。
考察其中一些推广将颇有收获。

最一般的拉氏量型 $\mathcal{N}=2$ QFT 有两类 $\mathcal{N}=2$ 超庞加莱代数的场表示：
矢量多重态与超多重态。

含矢量多重态与超多重态的理论同样可以拓扑扭曲；并且得到的理论同样落在第 3 节所述的
上同调拓扑场论的一般框架之内。

\subsubsection{超多重态}
\label{sec:4-7-1}

还有另一种场多重态——超多重态。

在 $\mathcal{N}=2$ SYM 理论中，可以再纳入一种场多重态，称为超多重态。规范群为 $G$ 的
理论中的一个超多重态由选取 $G$ 的一个四元数表示 $W$ 来定义。为简单起见（并以损失一些
一般性为代价），我们假定该表示可写成
\begin{equation}
  W=R\oplus R^{*},
  \tag{4.65}
\end{equation}
其中 $R$ 是 $G$ 的复表示。该场多重态可以看作在表示 $R\oplus R^{*}$ 中变换的一个
$\mathcal{N}=1$ 手征超场。一个 $\mathcal{N}=1$ 手征超场有一个领头分量，它就是复标量场，
于是我们得到在 $G$ 的 $R\oplus R^{*}$ 中变换的复标量 $q\oplus\tilde q$。在超多重态中，
若改而考虑
\begin{equation}
  M:=q\oplus\tilde q^{*},
  \tag{4.66}
\end{equation}
则得到在表示 $R$ 中变换的 $SU(2)_{R}$ 标量二重态。（注意 $SU(2)_{R}$ 不与超荷对易，
故 $M$ 不是 $\mathcal{N}=1$ 手征超场的底分量。）

现在，按上文所述对理论作拓扑扭曲时，$SU(2)_{R}$ 的二重态变成 $SU(2)'_{+}$ 的旋量。
因此，在 $X$ 上拓扑扭曲后的理论中，超多重态中的标量场成为截面
\begin{equation}
  M\in\Gamma(S^{+}\otimes\mathfrak{R}),\qquad
  \bar{M}\in\Gamma(S^{+}\otimes\mathfrak{R}^{*}),
  \tag{4.67}
\end{equation}
其中 $\mathfrak{R}\to X$ 现在是由表示 $R$ 与 $P\to X$ 相伴的向量丛。（超多重态中的
费米子在 $SU(2)_{R}$ 下不变，故也是自旋丛乘以 $\mathfrak{R}$ 的截面。）

这再次引出那个重要问题：必须能够表述 $X$ 非自旋、$w_{2}(X)\neq0$ 时的理论。解决办法与
把矢量多重态放到一般四流形上时相同：把 $S^{\pm}$ 视为“扭曲的”丛（扭曲 K-理论意义下），
再乘以“扭曲的”丛 $\mathfrak{R}$，使得 $S^{+}\otimes\mathfrak{R}$ 作为真正的丛存在。
因此必须取 $w_{2}(\mathfrak{R})=w_{2}(X)$，这或许要求我们为 $P$ 选取某个 't Hooft
磁通 $\xi$。

辅助场 $D$ 取值于 $\mathfrak{g}\otimes\mathbb{R}^{3}$，构成 $SU(2)_{R}$ 三重态。当我们
考虑超多重态处于四元数表示 $W$ 的矢量多重态 $+$ 超多重态拉氏量时，对 $D$-场作高斯消去
就把它等同于由 $G$ 在 $W$ 上的作用所定义的四元数矩映射
$\mu(M)\in\mathfrak{g}^{\vee}\otimes\mathbb{R}^{3}$。选取复结构使
$W\cong R\oplus\bar R$、$M=q\oplus\tilde q$，则四元数矩映射由二次表达式给出：
\begin{equation}
  \begin{aligned}
    T\cdot\mu_{r}(M)&=\langle q,T\cdot q\rangle
      -\langle\tilde q,T\cdot\tilde q\rangle,\\
    T\cdot\mu_{c}(M)&=\tilde q\cdot Tq .
  \end{aligned}
  \tag{4.68}
\end{equation}
左边我们用了 $\mathfrak{g}$ 与 $\mathfrak{g}^{\vee}$ 的自然配对，右边我们用了 $R$ 上的
厄米结构以及相应的同构 $\bar R\cong R^{\vee}$。由于 QFT 的作用量用到不变度量，可以把
$\mathfrak{g}\cong\mathfrak{g}^{\vee}$ 认同起来，而在对辅助场作高斯积分以置
$D=\mu$ 时用的正是这一认同。

于是可知，在拓扑扭曲理论中 $D$-场定义一个映射
\begin{equation}
  \mu:\Gamma(S^{+}\otimes\mathfrak{R})\to\Gamma(\Lambda^{+}\otimes\mathrm{ad}P),
  \tag{4.69}
\end{equation}
即从超多重态标量场到取值于 $\mathfrak{g}$ 的 $X$ 上自对偶二形式的映射。现在回忆
$\delta\chi=\ii(F^{+}-D^{+})$。连同超多重态中费米子的变分，$Q$ 不动点方程算出来是
\begin{equation}
  \begin{aligned}
    F^{+}&=\mu(M),\\
    \not{D}M&=0,
  \end{aligned}
  \tag{4.70}
\end{equation}
其中 $\not{D}:\Gamma(S^{+}\otimes\mathfrak{R})\to\Gamma(S^{-}\otimes\mathfrak{R})$ 是
“耦合到 $\mathfrak{R}$ 的自旋$^c$ 狄拉克算子”。这些方程称为\textbf{广义单极子方程}。
\footnote{〔原稿边注〕本节到此结束：距 §4.7 已过去 18 分钟。}

\subsubsection{$G=U(1)$ 理论中荷为 $1$ 的超多重态：Spin$^{c}$ 结构与塞伯格--威滕方程}
\label{sec:4-7-2}

广义单极子方程的一个特例在四流形（及三流形）理论中已被证明具有根本的重要性。我们取
规范群 $G=U(1)$、表示 $R\cong\mathbb{C}$ 为荷 $+1$ 的定义表示。

如上文所述，当 $X$ 非自旋时我们必须为 $S^{+}\otimes\mathfrak{R}$ 赋予意义；此时必须选取
一个自旋$^c$ 结构。物理学家的说法是：选取一个线丛 $L^{2}$，它只在局部才有平方根——
也就是说 $L$ 整体上并不存在；要求第一陈类满足
\begin{equation}
  w_{2}(X)=c_{1}(L^{2})\bmod 2,
  \tag{4.71}
\end{equation}
然后取 $M\in\Gamma(S^{+}\otimes L)$。由于过渡函数的上闭链（cocycle）关系中有 $-1$ 符号，
$S^{+}$ 与 $L$ 都不整体存在，但其乘积作为真正的向量丛存在。数学家通常把 $L^{2}$ 看作
一个秩二丛的行列式，即写作
\begin{equation}
  \mathrm{Det}(S^{+}\otimes L)\cong L^{2}.
  \tag{4.72}
\end{equation}
注意由于线丛 $L$ 未必存在，不能像普通线丛那样写 $c_{1}(L^{2})=2c_{1}(L)$，因此
$c_{1}(L^{2})\in H^{2}(X;\mathbb{Z})$ 未必能被 $2$ 整除！当它能被 $2$ 整除时，我们
可以找到一个自旋结构。

% [Tutor Note: (4.71)(4.72) 是自旋$^c$ 结构的“物理定义”：$L^2=\det(S^+\otimes L)$ 只
%  有整体定义，$L$ 本身至多在局部有定义，$c_1(L^2)$ 是 $w_2(X)$ 的整提升。(4.71) 中
%  的 mod 2 恰好是 $c_1(L^2)\equiv w_2(X)$ 这一同余，故 $c_1(L^2)$ 属于 $w_2(X)$ 的
%  整提升集合。(4.72) 里的 $\mathrm{Det}$ 是秩二丛的行列式线丛；它与 (4.75) 中的
%  $\lambda$、$\bar w_2(X)$ 的关系是 $2\lambda=c_1(L^2)$。]

我们刚刚给出的表述会让不少数学家感到不适。附录 A.3 给出了自旋$^c$ 结构在数学上更体面
的描述。

下面把这个情形下的广义单极子方程写出来。对 $U(1)$ 情形，用指标写法：若 $R$ 是 $U(1)$
的荷 $q$ 表示，则拓扑扭曲后 D 项是自对偶形式
\begin{equation}
  \mu(M)=q\bar{M}_{(\dot\alpha}M_{\dot\beta)},
  \tag{4.73}
\end{equation}
其中我们把自对偶形式与对称张量 $\Sym^{2}(W^{+})$ 认同起来。此时广义单极子方程就是与
自旋$^c$ 结构相联系的著名 \textbf{塞伯格--威滕方程}\footnote{它们有时也被称为
“单极子方程”。}：
\begin{equation}
  \begin{aligned}
    F^{+}_{\dot\alpha\dot\beta}
      &=\bar{M}_{(\dot\alpha}M_{\dot\beta)},\\
    D_{\dot\alpha\dot\beta}M^{\dot\beta}&=0 .
  \end{aligned}
  \tag{4.74}
\end{equation}
如我们已指出的，每个可定向四流形 $X$ 都是自旋$^c$ 的，因此塞伯格--威滕方程可以在
任意光滑的 $X$ 上写出。

\textbf{评注：}在下文 $u$-平面的计算中我们将遇到对自旋$^c$ 结构的求和，尤其会遇到与
自旋$^c$ 结构的挠子（torsor）相联系的 $\theta$ 函数。我们将有点马虎地忽略挠类，而直接把一个
自旋$^c$ 结构认同为元素
\begin{equation}
  \lambda\in\Gamma_{w}:=\frac{1}{2}\bar{w}_{2}(X)+\bar{H}^{2}(X),
  \tag{4.75}
\end{equation}
其中 $\bar{w}_{2}(X)$ 是 $w_{2}(X)$ 的一个固定整提升。按上述讨论的记号，
$2\lambda=c_{1}(L^{2})=c_{1}(\det W^{\pm})$。

\subsubsection{推广唐纳森--威滕配分函数}
\label{sec:4-7-3}

% [Tutor Note: (4.75) 的 $\Gamma_{w}$ 是一个“仿射格”：它是 $\bar H^2(X)$ 在
%  $\frac12\bar w_2(X)$ 处的平移，本身不是格而是挠子（torsor）。因此对 $\lambda$ 的
%  求和不是通常的格求和，而是带二次相位修正的 θ 函数求和：这正是 §6 中 $u$-平面积分
%  里出现带特征的 θ 函数（而非普通 θ 常数）的根源。作者特别声明忽略挠类，等价于只
%  考虑 $H^2(X;\mathbb{Z})$ 的自由部分。]

最一般经典 $\mathcal{N}=2$ 场论由以下数据定义：\footnote{〔原稿边注〕这里应当再多说一些，
解释一下克龙海默所做的工作。}
\begin{enumerate}[label=\arabic*., leftmargin=2.2em]
  \item 选取紧李群 $G$，并取其矢量多重态；
  \item 选取 $G$ 的四元数表示 $W$；
  \item 选取“质量参数”。
\end{enumerate}
于是可以写出经典作用量。该理论可以拓扑扭曲（受上文讨论的非自旋$^c$ 流形上丛的限制），
而拓扑扭曲后的理论其作用量落在第 3 节的上同调场论范式之内。\footnote{“质量参数”由
选取正交群中与规范群作用对易的子群的李代数的一个元素来定义。拓扑扭曲之后，这些参数可以
解释为整体对称性在 $Q$ 不动点模空间上的作用的等变参数。这一点在
[拉巴斯蒂达--马里诺]、\cite{ref26}\cite{ref24} 中有所讨论。}

鉴于上述评注，唐纳森理论似乎有一个巨大的推广。然而物理对数据加了很强的约束，只有这样
我们才能指望得到合理的答案：底层的物理 QFT 应当是良定义的。人们认为，只要所有规范耦合
的 $\beta$ 函数不是正的，该量子场论就是完全良定义的 QFT。若 $G$ 的李代数含有阿贝尔直和项，
这永远不会发生，故我们只应尝试对半单的 $G$ 定义推广。

若 $G$ 单且 $W\cong R\oplus\bar R$，则 $\beta$ 函数正比于
\begin{equation}
  \beta=-2h^{\vee}+C_{2}(R),
  \tag{4.76}
\end{equation}
其中 $C_{2}(R)$ 是 $R$ 的二次卡西米尔，归一化使 $C_{2}(\mathfrak{g}\otimes\mathbb{C})
=2h^{\vee}$。

最简单的推广是考虑单李群 $G$ 的纯矢量多重态理论。此时路径积分应局域化到群 $G$ 的 ASD
瞬子模空间上；由李代数上的卡西米尔还会生成更多可观测量。至少在形式层面，我们完全预期
威滕基本恒等式 (4.59) 的显然推广成立。

在数学一侧，这一推广似乎极其非平凡：ASD 模空间有额外的奇点；即使对无超多重态的
$SU(N)$，也不知道 Freed--Uhlenbeck“一般度量定理”的类比物。然而，\cite{ref22} 中
P. Kronheimer对一切 $N$ 给出了 $SU(N)$ 不变量的定义；关于高阶不变量更严格的数学处理
可见 \cite{ref5}。\footnote{当 $X$ 是代数曲面时，高阶不变量也已从代数几何的观点在
\cite{ref36} 中讨论过。}物理则提示它对任意单李群都可做到这一点。用下文所述 IR 方法对
$SU(N)$ 理论算出的“答案”已见于 \cite{ref31}。存在“威滕猜想”的推广，并且可以预期它
还能进一步推广到任意紧单李群，尽管这尚未被完成。
\footnote{〔原稿边注〕还是说已知它是错的……？}

更一般地，只要 $\beta$ 函数 (4.76) 非正，就应当能把超多重态包括进来。形式上，扭曲理论的
$Q$-闭可观测量的关联函数是在广义单极子方程 (4.70) 的解模空间上计算相交数。对
$G=SU(2)$、$R=\mathbf{2}$ 为 $SU(2)$ 基本表示的特殊情形，该模空间已由 P. Feehan与
莱内斯作了广泛讨论 \cite{ref14}；他们遇到的数学技术困难是令人生畏的。尽管如此，以 IR
理论表述的“答案”对 $G=SU(2)$、$R=N_{f}\mathbf{2}$（$N_{f}\le4$）确实已在
\cite{ref38} 中算出。\footnote{菲汉与莱内斯研究这个模空间，主要不是为了在它上面
做相交理论，而是把它用作 $SU(2)$ ASD 联络模空间与塞伯格--威滕模空间之间的配边。
这是为唐纳森不变量与塞伯格--威滕不变量之间关系给出严格数学证明可以采取的一条
途径。}把 UV 理论中具有非平凡超多重态表示的其它例子对应的 IR 理论算出来，或许是很有
启发的。我们将在讲义末尾回到这个问题以及物理所提示的其它推广。

% ---------------------------------------------------------------------
%  §4.8
% ---------------------------------------------------------------------
\subsection{那么，它有什么用？}
\label{sec:4-8}

威滕 1988 年的论文引入了拓扑场论的思想，尤其是拓扑扭曲的思想。这带来了对唐纳森多项式
的一个漂亮的量子场论诠释。

事后回看，这一诠释自然地提示塞伯格--威滕方程是 QFT 进路的必然产物，理由
很简单：把矢量多重态耦合到超多重态是自然的。事实上，确实有一些物理学家考虑过耦合超
多重态的拓扑扭曲理论。（见 \cite{ref24} 第 77 页及其所引文献 52,53,54。）\footnote{〔原稿
边注〕还有 Karlhede--Rocek。}广义单极子方程若作宽泛解释，可以说在这些文献中已被隐含地
定义，但并未被显式写出，方程式的几何内容也未被阐明。在威滕之前，似乎没有人意识到这些
方程的 $U(1)$ 版本的巨大威力。\footnote{〔原稿边注〕最后一句是否放错了位置？}

1988 年之后的若干年里，人们问道：“但这种诠释真的能导致计算唐纳森多项式的有效方法吗？”
这一点当时完全不清楚，若干怀疑论者持否定态度。确实存在相当程度的怀疑，甚至不无敌意，
直到 1994 年秋……\footnote{〔原稿边注〕本节到此结束：距 §4.7.2 已过去 25 分钟。}

% [Tutor Note: §4.8 的最后一句“直到 1994 年秋……”指的是塞伯格--威滕理论问世、
%  §5 以下要讲的红外方法给出唐纳森不变量显式公式的那一转折点。本节实际上把
%  “UV 定义 $\to$ IR 计算”这一全篇主线摆了出来：§4.6 的 $Z_{DW}^{\xi}(p,\Sigma)$ 由 UV
%  路径积分定义，而从 §5 起要在低能有效理论中重算同一个量。]
